
La démonstration du
théorème 3.6 occupera la fin de cet exposé. Donnons quelques préliminaires.

\textbf{Proposition 3.7.} \uline{Quels que soient} @@M0@@, @@M1@@, \uline{on a} :
@@M2@@
Par additivité, on peut supposer @@M3@@ et @@M4@@ homogènes; soit @@M5@@, @@M6@@.
Soient @@M7@@, @@M8@@ des éléments relevant @@M9@@ et @@M10@@. Alors la formule de projection donne :
@@M11@@
D'après 3.2, cette égalité s'écrit dans
@@M12@@. On en déduit une égalité dans
@@M13@@, et on vérifie immédiatement que c'est bien l'égalité cherchée.

\textbf{Lemme 3.8.} \uline{Soient} @@M14@@, \uline{et} @@M15@@ \uline{deux morphismes vérifiant les hypothèses de} 3.6. \uline{Alors} @@M16@@ \uline{les vérifie également}. \uline{Si} @@M17@@ \uline{est un élément de} @@M18@@, \uline{et si} 3.6 \uline{est vérifié par} @@M19@@ \uline{pour l'élément} @@M20@@, \uline{et par} @@M21@@ \uline{pour l'élément} @@M22@@, \uline{alors il est vérifié par} @@M23@@ \uline{pour l'élément} @@M24@@.

Par hypothèse, on a :
@@M25@@
et :
@@M26@@

En combinant (3.8.1) et (3.8.2), on obtient :
@@M27@@
